% Generated from hash-verified saved evidence. Do not edit by hand.
\begin{table}[htbp]
\centering
\caption{Illustrative posterior summaries from the first outer normal-prior fit.}
\label{tab:gen-posterior-examples}
\small
\setlength{\tabcolsep}{3pt}
\begin{adjustbox}{max width=\textwidth}
\begin{tabular}{llrrrr}
\toprule
Cohort & Feature & $E[\beta]$ & $E[\mathrm{HR}]$ & HR 95\% credible interval & $P(\beta>0)$ \\
\midrule
GBSG2 & Positive lymph nodes & 0.263 & 1.303 & [1.183, 1.420] & 1.000 \\
GBSG2 & Age & -0.122 & 0.890 & [0.721, 1.079] & 0.121 \\
GBSG2 & Progesterone receptor & -0.450 & 0.643 & [0.492, 0.808] & 0.000 \\
METABRIC & Age at diagnosis & 0.692 & 2.001 & [1.799, 2.217] & 1.000 \\
METABRIC & Nottingham prognostic index & 0.312 & 1.367 & [1.266, 1.471] & 1.000 \\
METABRIC & ER-positive & -0.359 & 0.704 & [0.539, 0.903] & 0.003 \\
WHAS500 & Age & 0.921 & 2.531 & [2.001, 3.181] & 1.000 \\
WHAS500 & Heart rate & 0.380 & 1.466 & [1.259, 1.692] & 1.000 \\
WHAS500 & Body mass index & -0.208 & 0.816 & [0.678, 0.971] & 0.011 \\
\bottomrule
\end{tabular}
\end{adjustbox}
\par\smallskip
\begin{minipage}{\textwidth}\footnotesize Examples use repeat 1, fold 1 and the normal prior, chosen by a fixed fold/prior rule rather than statistical significance. Continuous features are expressed per training-fold SD; categorical coefficients use the saved reference coding. HR is $\exp(\beta)$. These are model-conditional posterior associations, not pooled posteriors or causal effects.\end{minipage}
\end{table}
